% =====================================================================
%  Semantic Causal Propagation:
%  An Individuation-Theoretic Calculus of Boundary-Free Search,
%  and a Domain-Specific Language for Accountable Research Projects
%
%  Self-contained. Every object is a finite weighted graph or a
%  construction on one. No external framework or companion manuscript
%  is assumed; all citations are to standard, externally published
%  literature.
% =====================================================================
\documentclass[11pt,twocolumn,a4paper]{article}

\usepackage[utf8]{inputenc}
\usepackage[T1]{fontenc}
\usepackage{lmodern}
\usepackage[a4paper,margin=2.3cm]{geometry}
\usepackage{amsmath,amssymb,amsthm,mathtools}
\usepackage{bm}
\usepackage{mathrsfs}
\usepackage{booktabs}
\usepackage{array}
\usepackage{enumitem}
\usepackage{microtype}
\usepackage{graphicx}
\usepackage{xcolor}
\usepackage{listings}
\usepackage{algorithm}
\usepackage{algpseudocode}
\usepackage[round,sort&compress]{natbib}
\usepackage[colorlinks=true,linkcolor=blue!45!black,
            citecolor=blue!45!black,urlcolor=blue!45!black]{hyperref}
\usepackage{cleveref}

% ---- Theorem environments ----
\newtheoremstyle{thmstyle}{\topsep}{\topsep}{\itshape}{}{\bfseries}{.}{.5em}{}
\theoremstyle{thmstyle}
\newtheorem{theorem}{Theorem}[section]
\newtheorem{lemma}[theorem]{Lemma}
\newtheorem{proposition}[theorem]{Proposition}
\newtheorem{corollary}[theorem]{Corollary}
\theoremstyle{definition}
\newtheorem{definition}[theorem]{Definition}
\newtheorem{axiom}{Axiom}
\newtheorem{rulename}[theorem]{Rule}
\newtheorem{construction}[theorem]{Construction}
\newtheorem{example}[theorem]{Example}
\theoremstyle{remark}
\newtheorem{remark}[theorem]{Remark}
\newtheorem{notation}[theorem]{Notation}

\crefname{axiom}{Axiom}{Axioms}
\Crefname{axiom}{Axiom}{Axioms}
\crefname{rulename}{Rule}{Rules}
\Crefname{rulename}{Rule}{Rules}

% ---- Shorthand ----
\newcommand{\R}{\mathbb{R}}
\newcommand{\N}{\mathbb{N}}
\newcommand{\Z}{\mathbb{Z}}
\newcommand{\Gph}{\Gamma}
\newcommand{\V}{V}
\newcommand{\E}{E}
\newcommand{\wt}{w}
\newcommand{\med}{\mathsf{m}}
\newcommand{\flo}{\beta}
\newcommand{\cut}{\partial}
\newcommand{\co}{\complement}
\newcommand{\Walk}{W}
\newcommand{\rec}{M}
\newcommand{\res}{\rho}
\newcommand{\str}{\sigma}
\newcommand{\acc}{\mathcal{A}}
\newcommand{\Rep}{\mathsf{R}}
\newcommand{\Proc}{\Pi}
\newcommand{\pot}{\kappa}
\newcommand{\Cell}{\mathcal{C}}
\newcommand{\Cat}{\mathsf{Cat}}
\newcommand{\Proj}{\Xi}
\newcommand{\evalto}{\Downarrow}
\newcommand{\stepto}{\longrightarrow}
\newcommand{\bnfeq}{\;::=\;}
\newcommand{\bnfor}{\;\mid\;}
\newcommand{\abs}[1]{\lvert#1\rvert}
\newcommand{\norm}[1]{\lVert#1\rVert}
\newcommand{\eps}{\varepsilon}
\newcommand{\dd}{\,\mathrm{d}}
\newcommand{\nt}[1]{\ensuremath{\langle\textsf{#1}\rangle}}
\newcommand{\tm}[1]{\textbf{\texttt{#1}}}
\newcommand{\opt}[1]{\ensuremath{[\,#1\,]}}
\newcommand{\rep}[1]{\ensuremath{\{\,#1\,\}}}

% ---- grf listing style ----
\definecolor{grfkw}{RGB}{30,64,175}
\definecolor{grfcom}{RGB}{22,101,52}
\definecolor{grfnum}{RGB}{180,83,9}
\definecolor{grfstr}{RGB}{159,18,57}
\definecolor{grfbg}{RGB}{250,250,250}
\lstdefinelanguage{grf}{
  morekeywords={floor,project,seek,not,toward,via,until,converge,diverge,
    decline,yield,goal,catalyst,namespace,input,output,claims,coherence,
    let,if,then,else,import,module,export,assert,in,as,with},
  morekeywords=[2]{Claim,Region,Catalyst,Chain,Residue},
  sensitive=true,
  morecomment=[l]{--},
  morestring=[b]",
  keywordstyle=\color{grfkw}\bfseries,
  keywordstyle=[2]\color{grfnum},
  commentstyle=\color{grfcom}\itshape,
  stringstyle=\color{grfstr},
  basicstyle=\ttfamily\small,
  backgroundcolor=\color{grfbg},
  frame=single,framesep=4pt,xleftmargin=6pt,xrightmargin=4pt,
  breaklines=true,showstringspaces=false,columns=fullflexible,tabsize=2
}
\lstset{language=grf}

\title{\textbf{Semantic Causal Propagation:\\[3pt]
An Individuation-Theoretic Calculus of Boundary-Free Search,\\
and a Domain-Specific Language for Accountable Research Projects}}

\author{%
Kundai Farai Sachikonye\\
\small Technical University of Munich\\
\small AIMe Registry for Artificial Intelligence\\
\small \texttt{kundai.sachikonye@wzw.tum.de}}

\date{\today}

\begin{document}
\maketitle

% =====================================================================
\begin{abstract}
\noindent
We develop, entirely within the language of finite weighted graphs, a
formal theory of \emph{search} as the individuation of a claim from an
unbounded, non-completable medium, and a domain-specific orchestration
language, \emph{Graffiti} (\texttt{.grf}), whose executable semantics is
exactly this theory. The central observation is that a search query and
the act of naming or recognising the thing sought are the same operation
under two readings: both are the individuation of a region of claim-space
by negation against a medium that is never fully specified. From this we
derive, as theorems rather than assumptions: a strictly positive
\emph{resolution floor} on every search, forcing every returned claim to
be a bounded region rather than an exact point; a proof that natural-language
paraphrases of one underlying claim are \emph{representations} of a single
individuation and cost no additional search effort to substitute for one
another; a \emph{path-opacity} theorem guaranteeing that the interior
trajectory of a multi-step search is unconstrained provided the trajectory
converges on its target, so that intermediate results may look arbitrarily
unlike the sought claim without invalidating the search; a multiplicative
law for the combination of independent search steps (\emph{catalysts}),
together with a coherence theorem showing that no fewer than three mutually
reinforcing, independently sourced catalysts can ground a claim against
revision by a single dissenting source; and a \emph{closure} criterion for
knowing when a search is finished that is strictly stronger than crossing a
confidence threshold: a search is complete when every further catalyst
available to it resolves to a claim-region already reached, so that no
further inquiry can produce a materially different outcome. We then define
Graffiti as a small orchestration language whose sole computational
primitive, \texttt{seek}, is this individuation operation exposed to a
script author: a mandatory exclusion clause, a target claim-region, an
optional explicit catalyst chain, and a closure-based termination
condition. We give the complete lexical grammar, context-free grammar,
type system, and small-step operational semantics of Graffiti, prove
progress and preservation for its type system, prove that its scheduler
never revisits an already-closed claim-region and never starves a
convergent catalyst chain, and specify its execution model as a directed
acyclic graph of committed \texttt{seek} events over a heterogeneous
catalyst registry that may include local search, remote retrieval, and
machine-learned inference as computationally uniform operations. We
validate the framework on a suite of constructed instances and prove a
master equivalence tying the language's every guarantee to the single
fact that the medium of search is never completable.
\end{abstract}

\noindent\textbf{Keywords:} search theory; individuation by negation;
resolution floor; contact graphs; path opacity; representation mobility;
catalytic composition; coherence; domain-specific languages; operational
semantics; heterogeneous orchestration; accountable computation.

\tableofcontents

% =====================================================================
\part{Foundations}
% =====================================================================

% =====================================================================
\section{Introduction}
\label{sec:intro}
% =====================================================================

\subsection{The question}

What is a search? The question sounds settled by long practice --- one
supplies a query, a system returns candidates, and the practice of
information retrieval has an extensive engineering literature on ranking,
relevance, and recall~\citep{manning2008introduction,salton1983modern}.
But the engineering literature is silent on a structural question that
becomes acute the moment a search is not casual: when a user is not asking
for the nearest plausible page but is constructing a research artefact ---
a project, auditable and reproducible, whose steps matter as much as its
answer --- what, precisely, does the system return, when is it entitled to
stop, and what does it mean for two differently worded results to be
\emph{the same} finding?

We answer these questions from a single structural premise, argued rather
than assumed: a search is conducted against a \emph{medium} --- the totality
of what could bear on the query, without regard to whether it resides on
the querying machine or elsewhere --- and this medium is never exhausted by
any finite search process. From this one fact, formalised below as
non-completability (\Cref{def:noncompletable}), an entire calculus follows
by theorem: that no search can return an exact point rather than a bounded
region; that paraphrase costs nothing because it is not new search but a
change of coordinates on an already-found region; that a search's interior
steps may look arbitrarily unlike its answer without invalidating it,
provided the whole trajectory converges; that combining independent lines
of inquiry is a multiplicative operation with a precise law of diminishing
returns; and that a foundation-worthy claim requires no fewer than three
independently sourced, mutually reinforcing findings, never a single
authority nor even two.

\subsection{The recognition/search identity}

The premise that unifies the calculus is that \emph{recognising} a thing
and \emph{searching} for it are the same operation read in two directions.
To recognise an object is to place a percept into a region of meaning by
ruling out everything the percept is not; to search for an object is to
start from a description of what is wanted and to individuate, from the
medium, a region satisfying that description. Both operations terminate at
the same kind of object --- a bounded region of claim-space, never a single
point --- and both are subject to the same lower bound on how precisely
they can resolve their target. A query is not a degraded, lossy version of
a perfect specification the searcher secretly possesses; it is itself an
individuation, exactly as exact as recognition ever is. This identity,
established rigorously in \Cref{sec:individuation}, is what licenses
treating ``finding Peter,'' ``finding the founding date of an institution,''
and ``recognising that the object in front of you is a cat'' as instances
of one mathematical operation.

\subsection{Contributions}

\begin{enumerate}[leftmargin=1.4em,itemsep=2pt]
\item \textbf{The search medium and the resolution floor}
  (\S\ref{sec:medium}--\ref{sec:floor}). We formalise the medium of search
  as a non-completable expanding family of contact graphs and prove, from
  non-completability alone, that every search carries a strictly positive
  floor: no search returns a point, only a region.
\item \textbf{Individuation by negation and the identity of recognition and
  search} (\S\ref{sec:individuation}). We prove that in a medium with no
  privileged claim, a claim can be fixed only as the complement of what it
  is not, and that this is the same operation whether read as recognition
  (decoding a percept) or as search (projecting a description).
\item \textbf{Representation mobility} (\S\ref{sec:representation}). We
  prove that paraphrases of a fixed claim are equivalent representations
  under an averaging constraint, that switching between them commits no
  new search cost, and that no representation is privileged over another.
\item \textbf{Path opacity and convergence admissibility}
  (\S\ref{sec:path}). We prove that a multi-step search trajectory is
  admissible if and only if it converges on its target, with interior
  steps entirely unconstrained; two trajectories sharing endpoints are
  indistinguishable by any test of the output.
\item \textbf{Catalytic composition and the coherence theorem}
  (\S\ref{sec:catalysis}). We derive the multiplicative law for combining
  independent search steps, a Borel--Cantelli-type saturation dichotomy,
  and prove that a claim is robustly grounded if and only if its supporting
  catalysts form a strongly connected structure of at least three members.
\item \textbf{The closure criterion} (\S\ref{sec:closure}). We define
  \emph{closure} of a search --- the state in which no further catalyst
  resolves to a materially different claim-region --- and prove it is
  strictly stronger than, and the correct replacement for, a simple
  confidence threshold.
\item \textbf{The Graffiti language} (\S\ref{sec:language}--%
  \S\ref{sec:semantics}). We give the complete lexical grammar,
  context-free grammar, static type system, and small-step operational
  semantics of a language whose sole primitive, \texttt{seek}, is the
  individuation operation of \S\ref{sec:individuation}--\S\ref{sec:closure}
  exposed to a script author, and prove type soundness and scheduler
  soundness.
\item \textbf{The orchestration model} (\S\ref{sec:orchestration}). We
  specify Graffiti's execution model as a directed acyclic graph of
  committed \texttt{seek} events dispatched to a heterogeneous catalyst
  registry --- local computation, remote retrieval, and machine-learned
  inference treated as uniform operations --- and prove the scheduler is
  livelock-free and never re-executes a closed claim.
\end{enumerate}

\subsection{Method}

Every object in this paper is a finite weighted graph or a construction on
one, exactly as in the classical theory of network flows and cuts
\citep{fordfulkerson1956,mengertheorem,ahuja1993network}. We introduce no
metric space, no measure, and no probability calculus beyond a derived
normalisation on $[0,1]$. This is a discipline, not an aesthetic: every
claim reduces to a statement about vertices, edges, and weights, and is
therefore checkable by direct computation --- minimum cuts are exactly
computable by maximum flow~\citep{fordfulkerson1956,goldberg1988new}. Where
we use a word with an evocative reading --- ``claim,'' ``search,''
``catalyst'' --- the word names a specific graph-theoretic object defined
in the text, and no theorem depends on any reading beyond that definition.
Interpretive remarks are marked as such and are never load-bearing.

\subsection{Organisation}

\Cref{sec:medium} fixes the contact graph and the medium.
\Cref{sec:floor} derives the resolution floor from non-completability.
\Cref{sec:individuation} proves individuation by negation and the
recognition/search identity. \Cref{sec:representation} proves
representation mobility. \Cref{sec:path} proves path opacity and
convergence admissibility. \Cref{sec:catalysis} develops the catalytic
algebra and the coherence theorem. \Cref{sec:closure} defines and
characterises closure. \Cref{sec:language} gives the Graffiti lexical and
context-free grammar. \Cref{sec:types} gives the static type system and
proves soundness. \Cref{sec:semantics} gives the operational semantics and
proves determinism, monotonicity, and scheduler soundness.
\Cref{sec:orchestration} specifies the heterogeneous execution model.
\Cref{sec:examples} gives worked programs. \Cref{sec:validation} reports a
numerical validation suite. \Cref{sec:discussion} states scope and
concludes.

% =====================================================================
\section{The Search Medium}
\label{sec:medium}
% =====================================================================

\subsection{Contact graphs}

\begin{definition}[Finite weighted graph]
\label{def:graph}
A \emph{finite weighted graph} is a triple $\Gph=(\V,\E,\wt)$ where $\V$ is
a finite non-empty vertex set, $\E\subseteq\binom{\V}{2}$ is a set of
undirected edges, and $\wt\colon\E\to\R_{>0}$ is a strictly positive edge
weight. We assume $\Gph$ connected unless stated. Graphs are simple (no
loops or multi-edges).
\end{definition}

\begin{definition}[Contact graph; medium vertex]
\label{def:contact-graph}
A \emph{contact graph} is a finite weighted graph $\Gph=(\V,\E,\wt)$
together with a distinguished vertex $\med\in\V$, the \emph{medium},
adjacent to every other vertex: $\{v,\med\}\in\E$ for all $v\neq\med$.
Vertices $v\neq\med$ are \emph{claims}: candidate individuated units of
information --- a fact, a document, an entity, a numeric value, or any
other discrete thing a search may return. An edge $\{u,v\}$ between two
claims is a \emph{contact}, and its weight $\wt(\{u,v\})$ is the cost of
the distinction separating $u$ from $v$. The medium vertex $\med$ stands
for the undifferentiated totality of everything not yet individuated ---
every source, local or remote, that has not yet been resolved into a
claim.
\end{definition}

\begin{remark}[Why a single medium vertex, and why no boundary between
local and remote]
\label{rem:no-boundary}
The definition draws no distinction between a claim recoverable from a
file on the querying machine and a claim recoverable from a remote server:
both are, prior to being individuated, contents of the single vertex
$\med$. This is a deliberate structural commitment, not an approximation:
the theory that follows attaches no special status to physical locality,
because locality is not what makes a claim distinguishable from the rest
of the medium --- the cost of individuation (\Cref{def:sepcost}) is. Two
facts equally costly to establish are equally individuated regardless of
where the bytes that established them are stored. This licenses the design
principle of \Cref{sec:orchestration}: a script's catalysts may draw on
local files, remote APIs, or machine-learned models as computationally
uniform sources, because the graph never encodes a location, only a cost.
\end{remark}

\subsection{Separation cost}

\begin{definition}[Separation cost]
\label{def:sepcost}
The \emph{separation cost} of a claim $v$ is the minimum weight of a cut
placing $v$ on one side and the medium on the other,
\[
\str(v)\;=\;\min_{\substack{S\ni v\\ \med\notin S}}\wt(\cut(S)),
\qquad
\cut(S)\;:=\;\{\,e=\{u,w\}\in\E : u\in S,\ w\notin S\,\}.
\]
Because $\V$ is finite, the minimum is attained; the minimising side
$S^{\ast}(v)$ and the minimising cut $\cut(S^{\ast}(v))$ constitute the
\emph{minimum cut} of $v$ against the medium, computable exactly by
maximum flow~\citep{fordfulkerson1956,mengertheorem}.
\end{definition}

The separation cost is the graph-theoretic content of ``how hard it is to
tell $v$ apart from everything else.'' A claim recoverable from a single
uncontested primary source has low separation cost; a claim that can only
be established by triangulating many weakly corroborating traces has high
separation cost. Nothing in \Cref{def:sepcost} presumes what kind of
process realises a contact edge --- \Cref{sec:catalysis} shows the edges
themselves are supplied by \emph{catalysts}, and \Cref{sec:orchestration}
shows a catalyst may be a search engine query, a local file scan, or an
inference performed by a trained model.

% =====================================================================
\section{The Resolution Floor}
\label{sec:floor}
% =====================================================================

We now derive the central constant of the theory. It is not posited: it
follows from the conjunction of two facts already implicit in
\Cref{sec:medium} --- a claim is a proper part of the medium, and the
medium is never exhausted by any finite stage of search.

\subsection{Non-completability}

\begin{definition}[Expanding search family; non-completable medium]
\label{def:noncompletable}
An \emph{expanding search family} is a sequence of contact graphs
$\Gph_0\subseteq\Gph_1\subseteq\Gph_2\subseteq\cdots$, each obtained from
the previous by adding claims and contacts of weight $\ge\flo$ (fixed
below) while retaining the medium vertex and all prior claims and
contacts; $\Gph_k$ represents the state of search after $k$ committed
search steps. The medium is \emph{non-completable} if the family has no
terminal stage: for every finite $\Gph_k$ there exists $\Gph_{k+1}
\supsetneq\Gph_k$. This is an order condition --- no last stage of search
--- and not a cardinality assumption about how many claims exist.
\end{definition}

\begin{remark}[Why search is non-completable]
\label{rem:why-noncompletable}
A search medium spanning both a local machine and the unbounded remote
network admits, for any finite state of search, a further source not yet
consulted: another document, another database, another model that could be
queried, another cross-reference not yet followed. \Cref{def:noncompletable}
records only this order fact and nothing stronger; we make no claim about
the cardinality of the medium, only that no finite search process
exhausts it. This is the formal content of ``there is no boundary between
what is on the machine and what is on the internet'': the absence of a
boundary is exactly the absence of a terminal stage.
\end{remark}

\subsection{Identification and its residual}

\begin{definition}[Identification; identification residual]
\label{def:identification}
To \emph{identify} a claim $v$ at search stage $\Gph_k$ is to determine it
by its minimum cut against the medium, $\cut(S^{\ast}(v))$, computed within
$\Gph_k$. The identification is \emph{complete} if every claim of the whole
medium has entered the computation of that cut --- equivalently, if
$\Gph_k$ already contains the whole medium. The \emph{identification
residual} $\res(v)\ge0$ is the cut content a complete identification would
still remove that the finite-stage identification in $\Gph_k$ cannot;
$\res(v)=0$ if and only if the identification is complete.
\end{definition}

\subsection{The floor theorem}

\begin{theorem}[Floor from non-completability]
\label{thm:floor-derived}
Let the medium be non-completable and let $v$ be a claim ($v\neq\med$, and
$v$ does not exhaust $\V$). Then the identification of $v$ cannot be
completed at any finite search stage, and $\res(v)>0$.
\end{theorem}

\begin{proof}
Suppose $\res(v)=0$, so that identification of $v$ is complete at some
finite stage $\Gph_k$. By \Cref{def:identification}, completeness means
every claim of the medium has been brought to bear in the cut of $v$
against $\med$, i.e.\ the cut is computed in a graph containing the whole
medium; hence $\Gph_k$ contains the whole medium. This contradicts
non-completability (\Cref{def:noncompletable}), which guarantees a further
stage $\Gph_{k+1}\supsetneq\Gph_k$ containing a claim not present in
$\Gph_k$ and therefore not yet brought to bear on the cut of $v$.
Consequently no finite stage completes the identification of $v$, and
$\res(v)>0$.
\end{proof}

\begin{theorem}[Resolution floor]
\label{thm:floor}
There is a constant $\flo>0$ such that $\wt(e)\ge\flo$ for every contact
$e\in\E$ and $\str(v)\ge\flo$ for every claim $v$: no two claims are
separated at zero cost, and no claim is individuated for free.
\end{theorem}

\begin{proof}
Set $\flo:=\inf_v\res(v)$. By \Cref{thm:floor-derived}, $\res(v)>0$ for
every claim $v$; the residual is the irreducible boundary content of $v$'s
identification and is carried by the edges of $v$'s minimum cut, so each
such edge has weight at least the per-claim residual and hence at least
$\flo$. Because every contact edge $\{u,v\}$ is a boundary edge of the
identification of $u$ (it lies in $\cut(\{u\})$), every edge weight
satisfies $\wt(e)\ge\flo$. Connectivity of $\Gph$ forces
$\cut(S^{\ast}(v))\neq\varnothing$ for every claim $v$, so
$\str(v)=\wt(\cut(S^{\ast}(v)))\ge\flo>0$.
\end{proof}

\begin{corollary}[No sharp result]
\label{cor:no-sharp}
No search returns a claim separated from the medium by a cut of weight
zero. Every returned claim deposits boundary content $\ge\flo$; a search
that appears to answer ``exactly'' has, in fact, resolved to a region whose
width is bounded below by $\flo$ and no less.
\end{corollary}

\begin{remark}[The floor is a fact about the medium, not a defect of the
search process]
\label{rem:floor-fact}
\Cref{thm:floor} does not say search is imprecise because search engines
are imperfect; it says a search over a genuinely inexhaustible medium
cannot, even with an ideal instrument, terminate at a claim of zero
residual. A ``perfect'' search --- one returning a point rather than a
region --- would require completing a comparison against a medium that,
by hypothesis, is never completed. \Cref{sec:closure} makes this the basis
of an honest stopping rule rather than a caveat to be apologised for.
\end{remark}

% =====================================================================
\part{The Search Calculus}
% =====================================================================

% =====================================================================
\section{Individuation by Negation and the Recognition/Search Identity}
\label{sec:individuation}
% =====================================================================

\subsection{No privileged claim}

\begin{definition}[Selector; no privileged claim]
\label{def:selector}
A \emph{selector} for a claim set $U\subseteq\V$ is a rule that fixes $U$
by naming a distinguished member of $U$ in advance of any search. A
determination of $U$ is \emph{selector-free} if it names no member of $U$.
The medium has \emph{no privileged claim} if the data available to a
determining rule are invariant under the weighted automorphism group
$\operatorname{Aut}(\Gph)$, so that no claim is singled out by the data
alone.
\end{definition}

A medium with no privileged claim is the formal statement of a search
domain in which nothing is pre-marked as the answer: a claim must be
picked out by the structure of the query and the medium together, not by
an oracle that already knows which vertex is sought.

\subsection{Individuation by negation}

\begin{theorem}[Individuation by negation]
\label{thm:negation}
In a medium with no privileged claim, any selector-free determination of a
claim set $U$ fixes $U$ as the complement of its negation set,
$U=\V\setminus\co U$, with $\co U:=\V\setminus U$ specified first; and
conversely every admissible negation set $\co U$ determines a unique $U$.
No selector-based rule fixes $U$ without itself invoking a further
selector.
\end{theorem}

\begin{proof}
\emph{Necessity.} To fix $U$ is to decide, for every claim of $\V$, whether
it belongs to $U$ or to the remainder. A \emph{positive} rule --- one
asserting membership of a distinguished claim of $U$ --- must name that
claim; since the available data are $\operatorname{Aut}(\Gph)$-invariant
(\Cref{def:selector}), no claim is named by the data alone, so the rule
must itself supply a selector. That selector must in turn be fixed, either
given in advance (excluded by the no-privileged-claim hypothesis) or
supplied by a prior rule facing the identical dichotomy. The regress
terminates only at a rule that names no member of $U$: specifying the
negation set $\co U=\V\setminus U$ and taking $U$ as what remains. This
names only ``the rest,'' a datum already supplied once the medium $\V$ is
given.

\emph{Sufficiency and uniqueness.} Complementation $U\mapsto\V\setminus U$
is an involution on subsets of the finite set $\V$: $\V\setminus(\V
\setminus U)=U$. Hence a given negation set $\co U\subsetneq\V$ with
$\V\setminus\co U\neq\varnothing$ determines the unique claim set
$U:=\V\setminus\co U$, and conversely.
\end{proof}

\begin{corollary}[A claim is fixed by its non-neighbours]
\label{cor:claim-negation}
A single claim $v$ is determined selector-free by the partition of the
rest of the medium into its neighbours $N(v)$ (the claims it contacts) and
its non-neighbours $\V\setminus N[v]$ (the claims it neither is nor
contacts). No intrinsic label distinguishes $v$; its identity is carried
by this negation structure.
\end{corollary}

\begin{proof}
Apply \Cref{thm:negation} with $U=\{v\}$, refined by adjacency: the
selector-free datum fixing $\{v\}$ is $\V\setminus\{v\}$ partitioned into
$N(v)$ and $\V\setminus N[v]$; no member of $\{v\}$ itself is named by
this datum.
\end{proof}

\subsection{Decoding and projection: two readings of one map}

We now make precise the identity between recognising a claim (decoding a
given percept or query into its cell) and searching for a claim
(projecting a description onto the region of the medium that satisfies
it). We show these are literally the same function under an inverse
reading, not merely analogous procedures.

\begin{definition}[Query; decoder; projector]
\label{def:decoder}
Fix a contact graph $\Gph$. A \emph{query} is any finite string, structured
value, or percept $q$ presented to the search process. A \emph{decoder} is
a map $\mathrm{Dec}\colon Q\to\V$ from the space $Q$ of admissible queries
to claims of $\Gph$, sending a query to the claim it individuates by
\Cref{cor:claim-negation}. A \emph{projector} is a map
$\Proj\colon\V\to 2^{Q}\setminus\{\varnothing\}$ sending a claim to the
non-empty set of queries that decode to it: $\Proj(v):=\mathrm{Dec}^{-1}
(\{v\})$.
\end{definition}

\begin{theorem}[Recognition/search identity]
\label{thm:recognition-search}
For every claim $v\in\V$ and every query $q\in Q$:
\[
\mathrm{Dec}(q)=v \quad\Longleftrightarrow\quad q\in\Proj(v).
\]
Consequently \emph{recognition} --- computing $\mathrm{Dec}(q)$ for a given
percept $q$ --- and \emph{search} --- computing a representative element
of $\Proj(v)$ for a given target claim $v$, or equivalently testing
membership $q\in\Proj(v)$ for a candidate query $q$ --- are inverse
readings of the single relation
$\{(q,v)\in Q\times\V : \mathrm{Dec}(q)=v\}$. Neither is definable without
the other: $\Proj$ is literally the fibre structure of $\mathrm{Dec}$, and
$\mathrm{Dec}$ is the (well-defined, by \Cref{cor:claim-negation}) function
whose fibres are the sets $\Proj(v)$.
\end{theorem}

\begin{proof}
By \Cref{def:decoder}, $\Proj(v)=\mathrm{Dec}^{-1}(\{v\})$; membership
$q\in\Proj(v)$ is by definition of preimage equivalent to
$\mathrm{Dec}(q)=v$. This equivalence is definitional, not a coincidence to
be separately established, so the two directions state the same fact. The
claim that recognition and search are inverse readings of one relation
follows because a decoder is exactly a function and a projector is exactly
its fibre map; a function and its fibre map determine each other on a
finite domain: $\mathrm{Dec}$ is recovered from $\{\Proj(v)\}_{v\in\V}$ as
the unique map sending each $q$ to the $v$ with $q\in\Proj(v)$ (well
defined since the $\Proj(v)$ partition $Q$ under a total decoder), and
$\Proj$ is recovered from $\mathrm{Dec}$ by \Cref{def:decoder}.
\end{proof}

\begin{remark}[What this identity licenses]
\label{rem:identity-licenses}
\Cref{thm:recognition-search} is the formal content of the observation
that motivates this paper: a thought about a fixed object is the same
thought whether triggered by perceiving the object (recognition, computing
$\mathrm{Dec}$ forward) or by being asked to find it (search, computing a
representative of $\Proj$ backward). A search engine and a classifier are,
on this reading, the same computational object queried in opposite
directions. This licenses a single language primitive
(\Cref{sec:language}) for both operations: a \texttt{seek} expression names
a target region and asks the runtime to exhibit a query (a chain of
catalyst invocations) whose decode lands there, which is exactly computing
an element of $\Proj$ for that region.
\end{remark}

% =====================================================================
\section{Representation Mobility}
\label{sec:representation}
% =====================================================================

We now address paraphrase directly: many syntactically distinct
descriptions can individuate the same claim, and switching among them
must not be confused with performing additional search.

\subsection{Representations of a claim}

\begin{definition}[Representation; representation fibre]
\label{def:representation}
Fix a claim $v$ and a target claim $x^{\ast}$ against which alignment is
measured (\Cref{def:alignment} below formalises alignment; for now take
$a(v,x^{\ast})\in(0,1]$ as a fixed real). A \emph{representation} of $v$ at
dimension $N$ is a tuple $\Rep(v)=(s_1,\dots,s_N)\in\R^{N}$ satisfying the
\emph{averaging constraint}
\[
\frac{1}{N}\sum_{j=1}^{N} s_j \;=\; a(v,x^{\ast}),
\]
with individual components $s_j$ otherwise unconstrained in $\R$ (in
particular, a component may lie outside $(0,1]$). The \emph{representation
fibre} of $v$ at dimension $N$ is the set of all such tuples.
\end{definition}

\begin{theorem}[Representation mobility]
\label{thm:rep-mobility}
For every claim $v$ and every $N\ge2$:
\begin{enumerate}[label=\textup{(\roman*)},leftmargin=2.2em]
\item the representation fibre of $v$ at dimension $N$ is non-empty and
infinite;
\item switching from one representation of $v$ to another commits no new
contact edge and leaves the committed search-step count unchanged
(\Cref{def:committed} below);
\item any downstream computation that depends on $v$ only through the
average $a(v,x^{\ast})$ is invariant under representation switching.
\end{enumerate}
\end{theorem}

\begin{proof}
(i) The averaging constraint is a single linear equation in $N$ unknowns;
for $N\ge2$ its solution set is an affine hyperplane of $\R^{N}$, hence
non-empty and infinite, and any $(s_1,\dots,s_{N-1})\in\R^{N-1}$ extends
uniquely to a solution by $s_N=N\,a(v,x^{\ast})-\sum_{j<N}s_j$.

(ii) A representation switch replaces one tuple in $v$'s fibre with
another; both denote the same claim $v$ at the same alignment
$a(v,x^{\ast})$ (both satisfy the same averaging constraint), so the switch
introduces no new claim, no new contact, and hence no new committed search
step (\Cref{def:committed}).

(iii) By hypothesis the downstream computation is a function of
$a(v,x^{\ast})$ alone, and every member of $v$'s representation fibre
averages to the same value, so the computation's result is unchanged.
\end{proof}

\begin{example}[The stash example]
\label{ex:stash}
Let $v$ be a fixed individual, and let $x^{\ast}$ be the claim ``is
identifiable as this individual.'' The strings ``Peter,'' ``a Harvard
medical student,'' and ``a specialist in the anatomical processes of
diseased human physiology, training at an institution of higher learning
in Boston'' are three representations of $v$ in the sense of
\Cref{def:representation}: each, decoded by an appropriately informed
reader, individuates the same claim $v$ against the same target
$x^{\ast}$, and each satisfies the same averaging constraint even though
their surface content --- their component-wise detail --- differs
enormously. By \Cref{thm:rep-mobility}(ii), producing the third string
having already established the first commits no new search: it is a
change of representation on an already-individuated claim, not a
verification of a new one.
\end{example}

\begin{theorem}[Receiver-relative decoding is not error]
\label{thm:receiver-relative}
Let $\mathrm{Dec}_1,\mathrm{Dec}_2$ be two decoders (\Cref{def:decoder})
over graphs $\Gph_1,\Gph_2$ with distinct edge sets, modelling two readers
with different background knowledge. There exist a query $q$ (e.g.\ a
representation of $v$ as in \Cref{ex:stash}) and claims $v_1\neq v_2$ such
that $\mathrm{Dec}_1(q)=v_1$ and $\mathrm{Dec}_2(q)=v_2$, with both
decodings individuating a claim of separation cost $\ge\flo$ within their
respective graphs; neither decoding is in error.
\end{theorem}

\begin{proof}
Because $\Gph_1\neq\Gph_2$ as weighted graphs, the same query $q$ may
individuate different claims under $\mathrm{Dec}_1$ and $\mathrm{Dec}_2$
(the two decoders act on structurally different negation graphs, and
\Cref{thm:negation} determines $U$ relative to the ambient graph, not
absolutely). By \Cref{thm:floor} applied separately to $\Gph_1$ and
$\Gph_2$, both $v_1$ and $v_2$ satisfy $\str(v_i)\ge\flo_i>0$ in their
respective graphs, so both are legitimate individuations at their own
floor. Since \Cref{cor:claim-negation} makes each decoding a correct
application of individuation by negation within its own graph, neither
decoding fails any criterion the theory imposes; there is no privileged
$\Gph_i$ relative to which the other is judged incorrect.
\end{proof}

\begin{remark}[Why a ten-year-old and a physician are both right]
\label{rem:ten-year-old}
\Cref{thm:receiver-relative} is the formal statement that ``a Harvard
medical student'' registers a different, but equally legitimate, claim in
the decoder graph of a child (who may resolve only to the thin claim ``a
student'') than in the decoder graph of a colleague (who may resolve to a
rich claim with many further contacts: specialty, year of training,
research focus). Both are correct individuations at their own floor
(\Cref{thm:floor}); there is no third, privileged decoding of which both
are approximations. \Cref{sec:closure} uses this fact to justify why
closure of a search is always closure \emph{for a specified catalyst
registry and target audience}, never an absolute closure.
\end{remark}

% =====================================================================
\section{Path Opacity and Convergence Admissibility}
\label{sec:path}
% =====================================================================

We now make search dynamic: a search is not a single individuation but a
walk of successive contacts, and we characterise which walks are
admissible results.

\subsection{Alignment and propagation}

\begin{definition}[Alignment functional]
\label{def:alignment}
Fix a contact graph $\Gph$ and a target claim $x^{\ast}$. The
\emph{alignment} of a claim $x$ to $x^{\ast}$ is
\[
\str(x,x^{\ast})\;:=\;\min_{\substack{S\ni x\\ x^{\ast}\notin S}}
\wt(\cut(S))\;\in\;[\flo,\Omega],
\qquad \Omega:=\wt(\E),
\]
the minimum-cut weight separating $x$ from $x^{\ast}$, and the
\emph{alignment score} is $a(x,x^{\ast}):=\str(x,x^{\ast})/\Omega\in(0,1]$.
\end{definition}

\begin{lemma}[The floor bounds alignment below]
\label{lem:align-floor}
For distinct claims $x\neq x^{\ast}$, $a(x,x^{\ast})\ge\flo/\Omega>0$.
\end{lemma}

\begin{proof}
Immediate from \Cref{thm:floor}: any $x$--$x^{\ast}$ cut has weight
$\ge\flo$.
\end{proof}

\begin{definition}[Search process graph; propagation]
\label{def:process}
A \emph{search process graph} is a contact graph
$\Proc=(\V,\E,\wt;\,I,x^{\ast})$ together with a set $I\subseteq\V$ of
\emph{seed claims} (the initial information a search begins from --- a
query string decoded into an initial claim, a working hypothesis, or a
prior result) and a designated \emph{target claim} $x^{\ast}\in\V$. A
\emph{propagation} of a seed $v_0\in I$ through $\Proc$ is a walk
$\Walk=(v_0,e_1,v_1,\dots,e_k,v_k)$ along edges of $\Proc$ with
$v_k=x^{\ast}$; each edge $e_i$ is a \emph{contact operation}, realised in
\Cref{sec:catalysis} by a catalyst invocation.
\end{definition}

\begin{theorem}[Convergence admissibility]
\label{thm:admissibility}
A propagation $\Walk$ is \emph{admissible} if and only if it is
\emph{convergent}: its terminal claim is $x^{\ast}$. Interior claims
$v_1,\dots,v_{k-1}$ are entirely unconstrained --- an interior claim may
have arbitrarily large alignment $a(v_i,x^{\ast})$, i.e.\ may appear
maximally unrelated to the eventual target, without rendering $\Walk$
inadmissible, provided the walk terminates at $x^{\ast}$. Inadmissibility
arises exclusively from failure to reach $x^{\ast}$, never from local
misalignment along the way.
\end{theorem}

\begin{proof}
By \Cref{def:process}, $\Walk$ terminates at $x^{\ast}$ by construction
whenever it is a propagation in the sense defined; the terminal alignment
is then $a(x^{\ast},x^{\ast})=\flo/\Omega$, the least value permitted by
\Cref{lem:align-floor}, so every such walk is convergent. A walk failing to
reach $x^{\ast}$ has terminal claim $v_k\neq x^{\ast}$ and terminal
alignment $a(v_k,x^{\ast})\ge\flo/\Omega$, generally with strict
inequality, and is not a propagation to $x^{\ast}$ in the sense of
\Cref{def:process}; it is inadmissible as a search for $x^{\ast}$.
Nothing in \Cref{def:process} constrains the alignment of any interior
claim $v_i$, $0<i<k$, to the target: the walk is defined purely by
adjacency along $\Proc$ and by its terminus. Hence admissibility depends
exclusively on the terminal condition.
\end{proof}

\begin{theorem}[Path opacity]
\label{thm:path-opacity}
Let $\Walk_1,\Walk_2$ be two propagations of the same seed $v_0$ to the
same target $x^{\ast}$, differing in their interior claims. Then no
invariant computed from the seed and target alone --- the seed $v_0$, the
target $x^{\ast}$, the terminal alignment $a(x^{\ast},x^{\ast})$, or the
minimum cut $\cut(S^{\ast}(x^{\ast}))$ --- distinguishes $\Walk_1$ from
$\Walk_2$. The interior of a search trajectory is unobservable from its
result.
\end{theorem}

\begin{proof}
Every listed invariant is a function of $v_0$ and $x^{\ast}$ alone, hence
identical on $\Walk_1$ and $\Walk_2$ by hypothesis. The interior data ---
the sequence of intermediate claims and contacts --- are, by
\Cref{thm:admissibility}, unconstrained given the endpoints, so they are
not recoverable as a function of the endpoints. An invariant computed only
from the endpoints therefore cannot recover, and so cannot distinguish,
interior data that is not a function of those endpoints.
\end{proof}

\begin{corollary}[The admissible set is a class, not a path]
\label{cor:gauge}
For fixed $(v_0,x^{\ast},\Proc)$, the admissible propagations form an
equivalence class under ``shares seed and target, and converges,'' every
member of which is indistinguishable from every other by any test of the
result alone. A search's deliverable is properly this class, not a single
distinguished trajectory through it.
\end{corollary}

\begin{proof}
Immediate from \Cref{thm:admissibility,thm:path-opacity}: admissibility is
determined by the endpoints alone, and endpoint-invariant tests cannot
distinguish interiors, so all convergent same-endpoint propagations are
equivalent under any such test.
\end{proof}

\begin{remark}[A weird intermediate result does not invalidate a search]
\label{rem:weird-intermediate}
\Cref{thm:admissibility} licenses a concrete design consequence developed
in \Cref{sec:orchestration}: a multi-stage search script may pass through
an intermediate result that looks unrelated, contradictory, or simply
poor, and this is not by itself evidence the script has gone wrong. What
would be evidence of failure is the terminal claim's non-convergence ---
the script never reaching a claim within the target region --- never the
appearance of any single intermediate step.
\end{remark}

% =====================================================================
\section{Catalytic Composition and Coherence}
\label{sec:catalysis}
% =====================================================================

\subsection{Catalysts and catalytic power}

We now name the operation that realises a contact edge: the invocation of
a computational resource --- a search query, a local scan, a database
lookup, an inference from a trained model --- that moves a search closer to
its target.

\begin{definition}[Catalyst]
\label{def:catalyst}
A \emph{catalyst} toward target $x^{\ast}$ is a partial map
$\gamma\colon\V\to\V$ (acting on the current claim of a propagation) such
that, wherever defined, $\gamma$ corresponds to a contact operation of
\Cref{def:process}: applying $\gamma$ to a claim $x$ commits a contact
edge and yields $\gamma(x)$.
\end{definition}

\begin{definition}[Catalytic power]
\label{def:power}
The \emph{catalytic power} of a catalyst $\gamma$ toward $x^{\ast}$ at a
claim $x$ is the fraction of above-floor alignment it closes:
\[
\pot(\gamma,x,x^{\ast})\;=\;
\frac{a(x,x^{\ast})-a(\gamma(x),x^{\ast})}{a(x,x^{\ast})-\flo/\Omega}
\;\in\;[0,1].
\]
A catalyst is \emph{inert} ($\pot=0$) if it does not reduce alignment
toward $x^{\ast}$; it is \emph{absolute} ($\pot=1$) if it closes all
above-floor alignment in a single application.
\end{definition}

\begin{lemma}[Power is bounded]
\label{lem:powerbound}
For every catalyst, claim, and target, $\pot\in[0,1]$.
\end{lemma}

\begin{proof}
By \Cref{lem:align-floor}, $a(\gamma(x),x^{\ast})\ge\flo/\Omega$, so the
numerator of \Cref{def:power} never exceeds the denominator; non-negativity
holds for catalysts directed toward $x^{\ast}$ (a catalyst that increases
alignment is assigned $\pot=0$ by convention, since it is not a catalyst
toward $x^{\ast}$ in the operative sense).
\end{proof}

\subsection{The multiplicative law}

\begin{theorem}[Multiplicative composition]
\label{thm:multiplicative}
For a chain of catalysts $\gamma_1,\dots,\gamma_n$ applied in sequence
(each to the output of the last) with individual powers
$\pot_1,\dots,\pot_n$ toward a fixed target $x^{\ast}$, the composite power
is
\[
\pot(\gamma_1\catal\cdots\catal\gamma_n)
\;=\;1-\prod_{i=1}^{n}(1-\pot_i).
\]
\end{theorem}

\begin{proof}
Let $r_0=a(x_0,x^{\ast})-\flo/\Omega$ be the initial above-floor alignment
gap. Applying $\gamma_1$ leaves residual gap $r_1=(1-\pot_1)r_0$ by
\Cref{def:power}; applying $\gamma_2$ to the result leaves
$r_2=(1-\pot_2)r_1=(1-\pot_2)(1-\pot_1)r_0$; by induction on the chain
length, $r_n=\prod_{i=1}^{n}(1-\pot_i)\,r_0$. The fraction of the initial
gap closed by the whole chain is $1-r_n/r_0=1-\prod_i(1-\pot_i)$, which is
the composite power by \Cref{def:power} applied to the chain as a single
catalyst.
\end{proof}

\begin{corollary}[Diminishing returns]
\label{cor:diminishing}
Repeating a single catalyst of power $\pot$ a total of $n$ times yields
composite power $1-(1-\pot)^n$: strictly increasing in $n$, convergent to
$1$, but never reaching it for finite $n$. Each further repetition closes
a strictly smaller absolute fraction of the remaining gap than the last.
\end{corollary}

\begin{corollary}[Saturation dichotomy]
\label{cor:saturation}
A chain of catalysts with powers $(\pot_i)_{i\ge1}$ drives the residual
above-floor alignment gap to $0$ in the limit if and only if
$\sum_{i=1}^{\infty}\pot_i=\infty$.
\end{corollary}

\begin{proof}
The residual fraction after $n$ catalysts is $\prod_{i\le n}(1-\pot_i)$.
Taking logarithms, $\sum_{i\le n}\log(1-\pot_i)\to-\infty$ (residual
$\to 0$) if and only if $\sum_i\pot_i=\infty$, since $-\log(1-\pot)\asymp
\pot$ for small $\pot$ and the comparison test applies; this is a
Borel--Cantelli-type dichotomy~\citep{borel1909,cantelli1917,%
williams1991probability}.
\end{proof}

\begin{remark}[Design consequence: diversify catalysts, do not repeat
them]
\label{rem:diversify}
\Cref{cor:diminishing,cor:saturation} together justify a static check
performed by the Graffiti type checker (\Cref{sec:types}): a catalyst
chain consisting of near-duplicate invocations of one source (the same
search engine queried five paraphrased ways, say) accrues power slowly and
may fail to saturate even in principle, whereas a chain of independent
catalysts addressing distinct evidentiary channels reaches saturation
faster and more robustly. This is a structural fact about the
multiplicative law, not an empirical heuristic about search engines.
\end{remark}

\subsection{Coherence requires a triangle}

The multiplicative law of \Cref{thm:multiplicative} presumes every
catalyst in a chain is directed toward the same target. We now show that a
claim is robust against revision by a single dissenting source only if its
supporting catalysts form a cycle of at least three mutually reinforcing
members --- never a linear chain, and never merely two.

\begin{definition}[Support; support graph]
\label{def:support}
For catalysts $\gamma_i,\gamma_j$ used toward a common target $x^{\ast}$,
say $\gamma_j$ \emph{supports} $\gamma_i$ at threshold $\theta\in(0,1]$ if
the presence of $\gamma_j$'s result does not reduce, and to degree
$\ge\theta$ reinforces, the alignment shift toward $x^{\ast}$ induced by
$\gamma_i$. The \emph{support graph} $G$ of a set of catalysts has an
edge $j\to i$ whenever $\gamma_j$ supports $\gamma_i$ at $\theta$.
\end{definition}

\begin{theorem}[Linear support fails]
\label{thm:linear-fails}
No finite chain of catalysts $\gamma_1\to\gamma_2\to\cdots\to\gamma_n$, in
which each catalyst's contribution is justified only by the next and the
terminal catalyst is justified by nothing internal to the chain, can drive
the residual alignment gap to zero, nor can it be robust to removal of the
terminal catalyst.
\end{theorem}

\begin{proof}
The terminal catalyst $\gamma_n$ has no incoming support edge in $G$; its
contribution rests on nothing internal to the chain. By
\Cref{lem:align-floor}, no finite composite power reaches $1$ exactly
(\Cref{cor:diminishing}), so the residual gap is bounded below by the
floor regardless of chain construction; more materially, removing
$\gamma_n$ leaves $\gamma_{n-1}$ unsupported, and by induction each
successive removal leaves the next catalyst unsupported, so the chain's
claim to be grounded is contingent entirely on retaining every member: a
single removal at the terminus collapses the justification with no
internal check to prevent it.
\end{proof}

\begin{theorem}[Coherence requires three mutually supporting catalysts]
\label{thm:triangle}
\leavevmode
\begin{enumerate}[label=\textup{(\roman*)},leftmargin=1.8em,itemsep=2pt]
\item \emph{(Necessity.)} If a set of catalysts brings a search's terminal
alignment within any margin of the floor and remains so after the failure
of any single catalyst, its support graph contains a directed cycle of
length $\ge3$.
\item \emph{(Sufficiency.)} Three catalysts mutually supporting one
another at threshold $\theta>\tfrac12$, forming a strongly connected
triangle in $G$, ground the claim robustly: their combined shift toward
$x^{\ast}$ survives the removal of any single member.
\end{enumerate}
\end{theorem}

\begin{proof}
(i) By \Cref{thm:linear-fails}, no acyclic support structure is robust to
removal of its terminal member; an acyclic $G$ has a source catalyst with
no incoming support, and by the same argument its removal (or the removal
of any catalyst supporting only acyclically-downstream members) collapses
the chain below it. A $1$-cycle is self-support and vacuous (a catalyst
citing only itself). A $2$-cycle is mutual definition between exactly two
catalysts with no independent third check, and fails the majority
condition $\theta>\tfrac12$: two catalysts cannot outvote their own mutual
disagreement, since removing either leaves only one, itself unsupported.
Hence the shortest support cycle compatible with robustness has length
$\ge3$.

(ii) In a strongly connected triangle, each catalyst is supported by the
other two; with $\theta>\tfrac12$, the two supporters jointly exceed any
single dissenting signal. Removing one catalyst leaves the remaining two
still mutually supporting each other directly (the triangle's other edge),
above a reduced but still positive effective threshold, so the shift
toward $x^{\ast}$ persists.
\end{proof}

\begin{corollary}[Coherence is ordinally decidable]
\label{cor:decide}
Whether a set of catalysts is coherent in the sense of
\Cref{thm:triangle} is decidable from the \emph{signs} of pairwise support
relations alone --- whether each pair reinforces or undercuts the other's
shift toward $x^{\ast}$ --- with no access to catalytic-power magnitudes.
\end{corollary}

\begin{proof}
\Cref{thm:triangle} is stated and proved entirely in terms of the
directed-graph structure of $G$ (cycle existence and length) and the
threshold comparison $\theta>\tfrac12$, both of which are determined by
the signed support relation of \Cref{def:support} and not by the
numerical value of any catalytic power.
\end{proof}

% =====================================================================
\section{Closure: When a Search Is Finished}
\label{sec:closure}
% =====================================================================

We now state the paper's central operational contribution: a criterion for
terminating a search that is strictly stronger than ``confidence exceeds a
threshold,'' and is instead a structural fact about the space of further
catalysts available.

\subsection{Closure of the admissible set}

\begin{definition}[Reachable target-cell set]
\label{def:reachable}
Fix a seed $v_0$ and a search process graph $\Proc$ whose available
catalysts (edges) may be extended over time as new catalysts are
registered or invoked (\Cref{sec:orchestration}). At a given search stage,
let $\Cell(v_0,\Proc)$ denote the set of target claims $x^{\ast}$ reachable
by some admissible propagation of $v_0$ through the catalysts invoked so
far (\Cref{def:process,thm:admissibility}), partitioned into equivalence
classes under ``endpoint-indistinguishable'' (\Cref{cor:gauge}).
\end{definition}

\begin{definition}[Closure]
\label{def:closure}
A search of seed $v_0$ is \emph{closed} at a given stage if, for every
catalyst $\gamma$ available but not yet invoked, extending the search by
$\gamma$ cannot add a new equivalence class to $\Cell(v_0,\Proc)$: every
further admissible propagation reachable via $\gamma$ terminates in a
class already present in $\Cell(v_0,\Proc)$.
\end{definition}

\begin{theorem}[Closure is strictly stronger than a confidence threshold]
\label{thm:closure-stronger}
There exist search process graphs and confidence thresholds $\theta$ such
that a propagation's terminal alignment satisfies $a(x^{\ast},x^{\ast})
\ge\theta$ (a fixed confidence criterion is met) while the search is not
closed in the sense of \Cref{def:closure}: a further, not-yet-invoked
catalyst is available whose propagation reaches a target in a distinct
equivalence class.
\end{theorem}

\begin{proof}
Construct $\Proc$ with two disjoint claim clusters $A,B$, each internally
dense (many low-cost internal contacts) and each connected to the medium
by a single catalyst, $\gamma_A$ and $\gamma_B$ respectively, of equal
weight. A propagation via $\gamma_A$ alone terminates at a representative
claim $x^{\ast}_A\in A$ with alignment $a(x^{\ast}_A,x^{\ast}_A)=\flo/\Omega$,
satisfying any fixed threshold $\theta\le1$ trivially, since a
propagation's terminal alignment to \emph{itself} is always at the floor
by \Cref{thm:admissibility}. Yet $\gamma_B$, not yet invoked, reaches
$x^{\ast}_B\in B$, a claim in a distinct equivalence class from
$x^{\ast}_A$ (the two clusters are, by construction, not
endpoint-indistinguishable, since e.g.\ their minimum cuts against the
medium differ if $A\neq B$ as vertex sets with independent internal
structure). Hence the confidence criterion is satisfied after $\gamma_A$
alone, but the search is not closed: $\gamma_B$ can still add a new
equivalence class.
\end{proof}

\begin{remark}[Why closure is the correct criterion]
\label{rem:closure-correct}
\Cref{thm:closure-stronger} formalises a common failure of confidence-
threshold stopping rules: a single early source may report high confidence
purely because it is internally self-consistent (as any single cluster
$A$ is, having only committed contacts within itself and one to the
medium), while a second, independent source would reach an entirely
different conclusion. Closure requires that \emph{every} available further
catalyst be checked to land in an already-reached equivalence class before
stopping, which is exactly the condition ``no other possible paths can
still be generated'' motivating this definition: it is not merely that the
current path looks confident, but that the space of paths has stopped
producing new destinations.
\end{remark}

\subsection{Two ways a search can end}

\begin{theorem}[Convergent closure or honest decline]
\label{thm:decline}
Every search of seed $v_0$ over a search process graph with finitely many
registered catalysts terminates in exactly one of two states:
\begin{enumerate}[label=\textup{(\roman*)},leftmargin=2.2em]
\item \textbf{Convergent closure}: $\Cell(v_0,\Proc)$ stabilises to a
single equivalence class, and the search reports that class's
representative as its result;
\item \textbf{Contested closure}: $\Cell(v_0,\Proc)$ stabilises (no further
catalyst adds a class) but contains more than one equivalence class, and
the search reports \emph{decline}: no single sufficient claim exists,
together with the distinct classes found.
\end{enumerate}
No search over a finite catalyst registry runs forever without reaching
one of these two states.
\end{theorem}

\begin{proof}
Because the catalyst registry available to a given search is finite at any
call (\Cref{sec:orchestration} fixes registries as finite sets), the set
of propagations reachable from $v_0$ is finite, and hence the set of
equivalence classes $\Cell(v_0,\Proc)$ --- being a partition of a finite
set of reachable targets --- is finite and can only grow by finitely many
steps before every remaining catalyst has been tried against every
currently-reached class. At that point \Cref{def:closure} is satisfied by
exhaustion, so closure is reached in finitely many search steps. If the
stabilised $\Cell(v_0,\Proc)$ has exactly one class, the search is
convergently closed (i); if it has more than one, no single sufficient
claim is licensed by \Cref{cor:gauge} applied separately to each class
(each is internally endpoint-indistinguishable, but the classes disagree
with one another), and the search reports decline (ii). These two cases
are exhaustive and mutually exclusive on a partition with at least one
class (the search graph is non-empty by hypothesis, so at least one class
exists once any propagation is admissible).
\end{proof}

\begin{remark}[Decline is not failure]
\label{rem:decline-not-failure}
\Cref{thm:decline}(ii) formalises the requirement, motivated independently
in \Cref{sec:orchestration}, that a search runtime must be able to report
``the sources disagree'' as a first-class, correctly-terminating outcome,
rather than either looping indefinitely or silently emitting one of the
contested classes as if it were unique. Contested closure is itself useful
research output: it identifies precisely where independent inquiry
diverges.
\end{remark}

% =====================================================================
\part{The Graffiti Language}
% =====================================================================

% =====================================================================
\section{Lexical Structure and Grammar}
\label{sec:language}
% =====================================================================

We now specify \emph{Graffiti} (file extension \texttt{.grf}), a language
whose sole computational primitive, \texttt{seek}, is exactly the search
process of \Cref{sec:path,sec:catalysis,sec:closure} exposed to a script
author, and whose top-level unit of organisation, \texttt{project}, is a
directed acyclic graph of committed \texttt{seek} events.

\subsection{Design principles}

\begin{enumerate}[leftmargin=1.4em,itemsep=2pt]
\item \textbf{One primitive, four clauses.} Every \texttt{seek} names a
mandatory exclusion (\texttt{not}, \Cref{thm:negation}), a target region
(\texttt{toward}, \Cref{def:process}), an optional explicit catalyst chain
(\texttt{via}, \Cref{sec:catalysis}), and a termination condition
(\texttt{until}, \Cref{sec:closure}).
\item \textbf{Exclusion is mandatory.} A \texttt{seek} with no
\texttt{not} clause is rejected at parse time: by \Cref{thm:negation}, a
claim in a medium with no privileged claim is individuated only by what it
is not, so an expression asserting no exclusion asserts no individuation.
\item \textbf{Projects are graphs, not scripts.} A \texttt{project} is a
finite set of named \texttt{seek} bindings that may reference each other's
results, forming a directed acyclic graph; later bindings narrow using
the exclusions and yields of earlier ones.
\item \textbf{Closure, not confidence, is the default stopping rule.}
\texttt{until converge} means closure in the sense of
\Cref{def:closure}, not merely crossing a numeric threshold; a script may
opt into a threshold explicitly if a weaker guarantee is acceptable.
\item \textbf{Decline is a value, not an exception.} A \texttt{seek} that
reaches contested closure (\Cref{thm:decline}(ii)) yields a
\texttt{decline} value carrying the distinct equivalence classes found,
which downstream bindings may inspect.
\end{enumerate}

\subsection{Lexical grammar}

\begin{definition}[Token classes]
\label{def:tokens}
A Graffiti source file is UTF-8 text with extension \texttt{.grf}.
Lexing produces a token stream from the classes:
\begin{itemize}[leftmargin=1.4em,itemsep=1pt]
\item \textbf{Comments.} \texttt{--} to end of line; discarded.
\item \textbf{Keywords.} \texttt{floor}, \texttt{project}, \texttt{seek},
\texttt{not}, \texttt{toward}, \texttt{via}, \texttt{until},
\texttt{converge}, \texttt{diverge}, \texttt{decline}, \texttt{yield},
\texttt{goal}, \texttt{catalyst}, \texttt{namespace}, \texttt{input},
\texttt{output}, \texttt{claims}, \texttt{coherence}, \texttt{let},
\texttt{if}, \texttt{then}, \texttt{else}, \texttt{import},
\texttt{module}, \texttt{export}, \texttt{assert}, \texttt{in},
\texttt{as}, \texttt{with}.
\item \textbf{Identifiers.} \texttt{[A-Za-z\_][A-Za-z0-9\_]*}, excluding
keywords.
\item \textbf{Numbers.} decimal integers and floats with optional
exponent: \texttt{[0-9]+(\allowbreak.[0-9]+)?\allowbreak([eE][+-]?[0-9]+)?}.
\item \textbf{Strings.} double-quoted, with standard escapes.
\item \textbf{Operators and punctuation.} \texttt{:=} (bind),
\texttt{>>} (catalyse, sequential composition of catalysts),
\texttt{\textbar\textbar} (parallel catalyst dispatch),
\texttt{( ) \{ \} [ ] , : . > < >= <= == ->} and the field-access
operator \texttt{.}.
\end{itemize}
Whitespace is insignificant except as a token separator; Graffiti is not
layout-sensitive.
\end{definition}

\subsection{Context-free grammar}

We give the grammar in extended Backus--Naur
form~\citep{backus1959,naur1963,aho2006}. Nonterminals are
\textit{italicised}; terminals are \texttt{typewriter}; \texttt{\{x\}} is
zero-or-more repetition; \texttt{[x]} is optional; \texttt{|} is
alternation.

\begin{definition}[Graffiti grammar]
\label{def:grammar}
\begin{align*}
\textit{program} &\bnfeq \{\textit{decl}\} \\[2pt]
\textit{decl} &\bnfeq \textit{floorDecl} \bnfor \textit{import}
   \bnfor \textit{moduleDecl} \bnfor \textit{projectDecl}
   \bnfor \textit{catalystDecl}\\[2pt]
\textit{floorDecl} &\bnfeq \tm{floor}\ \textit{number} \\
\textit{import} &\bnfeq \tm{import}\ \textit{qname} \\
\textit{moduleDecl} &\bnfeq \tm{module}\ \textit{ident}\ \tm{\{}\
   \{\textit{decl}\}\ \tm{\}} \\[2pt]
\textit{catalystDecl} &\bnfeq \tm{catalyst}\ \textit{ident}\ \tm{\{}\\
   &\quad \tm{namespace}\tm{:}\ \textit{ident}\ \\
   &\quad \tm{input}\tm{:}\ \textit{type}\ \tm{output}\tm{:}\ \textit{type}\
   \tm{\}}\\[4pt]
\textit{projectDecl} &\bnfeq \tm{project}\ \textit{ident}\ \tm{\{}\
   \{\textit{seekStmt}\}\\
   &\quad [\textit{goalBlock}]\ \tm{\}} \\[4pt]
\textit{seekStmt} &\bnfeq \tm{seek}\ \textit{ident}\\
   &\quad \tm{not}\ \tm{\{}\ \textit{boundary}\ \tm{\}}\\
   &\quad \tm{toward}\ \tm{\{}\ \textit{region}\ \tm{\}}\\
   &\quad [\,\tm{via}\ \tm{\{}\ \textit{chain}\ \tm{\}}\,]\\
   &\quad \tm{until}\ \textit{admit}\\
   &\quad \tm{yield}\ \textit{ident} \\[4pt]
\textit{admit} &\bnfeq \tm{converge}\ [\,\tm{otherwise}\ \tm{decline}\,]\\
   &\bnfor \textit{cond} \\[4pt]
\textit{chain} &\bnfeq \textit{catalystRef}\ \{\tm{>>}\ \textit{catalystRef}\}\\
   &\bnfor \textit{chain}\ \tm{\textbar\textbar}\ \textit{chain} \\
\textit{catalystRef} &\bnfeq \textit{qname}\ \tm{(}\ [\textit{argList}]\ \tm{)} \\[4pt]
\textit{boundary} &\bnfeq \textit{ident}\ \{\tm{,}\ \textit{ident}\}\\
   &\bnfor \textit{boundary}\ \tm{\textbar\textbar}\ \textit{boundary}\\
   &\bnfor \textit{boundary}\ \tm{\&\&}\ \textit{boundary} \\[4pt]
\textit{region} &\bnfeq \textit{qname}\ \tm{(}\ [\textit{argList}]\ \tm{)}\\
   &\bnfor \textit{ident} \\[4pt]
\textit{goalBlock} &\bnfeq \tm{goal}\ \tm{\{}\\
   &\quad \tm{claims}\tm{:}\ \tm{[}\ \textit{ident}\
     \{\tm{,}\ \textit{ident}\}\ \tm{]}\\
   &\quad \tm{coherence}\tm{:}\ \textit{relop}\ \textit{number}\
   \tm{\}} \\[4pt]
\textit{argList} &\bnfeq \textit{arg}\ \{\tm{,}\ \textit{arg}\} \\
\textit{arg} &\bnfeq [\textit{ident}\ \tm{:}]\ \textit{expr} \\
\textit{cond} &\bnfeq \textit{expr}\ \textit{relop}\ \textit{expr} \\
\textit{relop} &\bnfeq \tm{>} \bnfor \tm{<} \bnfor \tm{>=} \bnfor \tm{<=}
   \bnfor \tm{==} \\
\textit{qname} &\bnfeq \textit{ident}\ \{\tm{.}\ \textit{ident}\} \\
\textit{expr} &\bnfeq \textit{ident} \bnfor \textit{number} \bnfor
   \textit{string} \bnfor \textit{catalystRef} \bnfor \tm{(}\ \textit{expr}\
   \tm{)}\\
\textit{type} &\bnfeq \tm{Claim} \bnfor \tm{Region} \bnfor \tm{Catalyst}
   \bnfor \tm{Chain} \bnfor \tm{Residue}
\end{align*}
\end{definition}

\begin{remark}[Reading the core forms]
\label{rem:core-forms}
A \texttt{seek} statement is the syntactic image of
\Cref{def:process,thm:admissibility,def:closure}: \texttt{not\{\}} supplies
the negation set of \Cref{thm:negation}; \texttt{toward\{\}} names the
target claim $x^{\ast}$; \texttt{via\{\}}, when present, supplies an
explicit catalyst chain (\Cref{sec:catalysis}) rather than leaving chain
derivation to the runtime; \texttt{until converge} invokes closure
(\Cref{def:closure}), and the optional \texttt{otherwise decline} clause
names the contested-closure outcome of \Cref{thm:decline}(ii) explicitly.
A \texttt{project} is a \textit{projectDecl}: an ordered set of
\textit{seekStmt}s whose \texttt{yield} identifiers may be referenced as
arguments to later catalyst invocations, giving the directed acyclic
dependency structure specified operationally in \Cref{sec:orchestration}.
\end{remark}

% =====================================================================
\section{The Type System}
\label{sec:types}
% =====================================================================

Graffiti's type system exists to make two properties of
\Cref{sec:individuation,sec:catalysis} statically checkable: that no claim
of zero residue is ever assigned a type (an accountability property, in
the sense that every value must be traceable to a floor-respecting
individuation), and that a coherence-requiring project cannot type-check
with fewer than the three mutually supporting catalysts
\Cref{thm:triangle} proves necessary.

\subsection{Types}

\begin{definition}[Types]
\label{def:types}
The Graffiti types are
\[
\tau \bnfeq \tm{Claim} \bnfor \tm{Region} \bnfor \tm{Catalyst}
   \bnfor \tm{Chain} \bnfor \tm{Residue}.
\]
\texttt{Claim} is the type of an individuated result (a \texttt{seek}
yield); \texttt{Region} is the type of a \texttt{toward\{\}} target,
before individuation; \texttt{Catalyst} is the type of a single named
computational resource (\Cref{sec:orchestration}); \texttt{Chain} is the
type of a composed catalyst sequence; \texttt{Residue} is the type of a
measured alignment gap (\Cref{def:alignment}). Every value of type
\texttt{Claim} additionally carries a \emph{floor annotation}
$\flo(v)>0$ and a \emph{residue} $\res(v)\ge\flo(v)$, exactly as in
\Cref{def:residue}.
\end{definition}

\begin{definition}[Typing judgement]
\label{def:judgement}
A judgement $\Gamma\vdash e:\tau\,@\,\flo$ reads: ``in floor context
$\Gamma$ (fixing the ambient floor and the types of bound identifiers),
expression $e$ has type $\tau$ and, if it individuates a claim, does so at
floor $\flo>0$.''
\end{definition}

\subsection{The positivity rule}

\begin{rulename}[Floor positivity]
\label{rule:positivity}
Every claim-introducing rule carries the side condition $\flo>0$ and
$\res\ge\flo$:
\[
\frac{\Gamma\vdash e_1:\tm{Region}\quad
      \Gamma\vdash e_2:\tm{Chain}\quad \flo_{\Gamma}>0}
     {\Gamma\vdash \tm{seek}\ n\ \tm{not}\{\cdot\}\ \tm{toward}\{e_1\}\
      \tm{via}\{e_2\}\ \cdots : \tm{Claim}\,@\,\flo_{\Gamma}}
\]
where $\flo_{\Gamma}$ is the ambient floor set by the most recent
\texttt{floor} declaration in scope.
\end{rulename}

\begin{theorem}[No zero-residue claim]
\label{thm:no-zero}
There is no well-typed Graffiti expression of residue $0$: every
typeable \texttt{Claim} value has $\res\ge\flo>0$.
\end{theorem}

\begin{proof}
By induction on typing derivations. The only rule introducing a
\texttt{Claim} value is the \texttt{seek} rule (\Cref{rule:positivity}),
whose side condition requires $\flo_{\Gamma}>0$; by \Cref{thm:floor} the
residue of the individuated claim satisfies $\res\ge\flo_{\Gamma}$
whenever the ambient floor is positive. No other rule introduces a
\texttt{Claim} value, and no rule can lower the residue annotation of an
already-typed value (representation switching, by
\Cref{thm:rep-mobility}(ii), preserves the alignment and hence the
residue exactly). Hence no derivation concludes $\res=0$.
\end{proof}

\subsection{The coherence rule}

\begin{rulename}[Coherence for closure]
\label{rule:coherence}
A \texttt{seek} statement whose \texttt{via\{\}} clause supplies an
explicit chain, and whose \texttt{until} clause is \texttt{converge}
(without an explicit \texttt{cond} weakening it), type-checks only if the
supplied chain's induced support graph (\Cref{def:support}) contains a
strongly connected subgraph on at least three catalysts:
\[
\frac{\Gamma\vdash e_2:\tm{Chain}\quad
      \text{supportGraph}(e_2)\text{ has a $\ge3$-cycle}}
     {\Gamma\vdash \tm{seek}\ n\ \cdots\ \tm{via}\{e_2\}\ \tm{until}\
      \tm{converge}\ \cdots : \tm{Claim}\,@\,\flo_{\Gamma}}
\]
A chain failing the cycle condition emits a static
\texttt{CoherenceWarning} rather than a hard type error, since a
\texttt{seek} lacking a triangle may still be legitimate under an
explicit, weaker \texttt{cond} admission criterion (a script author may
deliberately accept single-source or two-source evidence, provided this is
not silently mistaken for full closure).
\end{rulename}

\begin{theorem}[Type soundness]
\label{thm:soundness}
Graffiti typing is sound for the operational semantics of
\Cref{sec:semantics}: if $\Gamma\vdash e:\tau\,@\,\flo$ and $e$ evaluates
to a value $v$, then $v$ has type $\tau$, floor $\flo>0$, and (if
$\tau=\tm{Claim}$) residue $\res(v)\ge\flo$.
\end{theorem}

\begin{proof}[Proof sketch]
Standard syntactic type soundness by progress and
preservation~\citep{pierce2002types,wright1994syntactic}. \emph{Progress}:
a well-typed closed expression is either a value or takes a step, by
inspection of the reduction rules of \Cref{sec:semantics}, each of which
has a syntactically disjoint premise from every other rule of the same
constructor. \emph{Preservation}: each reduction rule preserves the
static type and floor annotation of \Cref{rule:positivity}, and, by
\Cref{thm:no-zero} applied at every step, preserves $\res\ge\flo>0$ for
every \texttt{Claim} value produced. The only non-standard element is that
reduction is state-mutating (\Cref{sec:semantics}: evaluating a
\texttt{seek} commits contact edges and advances the search-step count);
preservation nonetheless holds because the committed-step count is
monotone (\Cref{thm:monotone-count} below), so no reduction step can
retroactively lower a previously-established residue below its floor.
\end{proof}

% =====================================================================
\section{Operational Semantics}
\label{sec:semantics}
% =====================================================================

We give a small-step semantics over configurations
$\langle\Gph,\rec,e\rangle$, where $\Gph$ is the current contact graph
(the accumulated state of search so far), $\rec\in\N$ is the committed
search-step count, and $e$ is the expression under evaluation. We write
$\langle\Gph,\rec,e\rangle\stepto\langle\Gph',\rec',e'\rangle$ for one
step. Evaluation is state-mutating: a \texttt{seek} step commits contact
edges into $\Gph$ and advances $\rec$.

\subsection{Committed steps and the intrinsic clock}

\begin{definition}[Committed step; committed record]
\label{def:committed}
A \emph{committed step} is a reduction that commits one contact edge into
the current graph $\Gph$ (a single catalyst invocation resolving to a new
claim or refining an existing one). The \emph{committed record}
$\rec\in\N$ is the count of committed steps taken so far in the current
evaluation.
\end{definition}

\begin{theorem}[Monotonicity of the committed record]
\label{thm:monotone-count}
Along any evaluation, $\rec$ is non-decreasing, and strictly increases at
every committed step. No reduction rule decrements $\rec$; re-evaluating a
\texttt{seek} expression already bound is a new committed step at a
strictly higher $\rec$, never a cached recomputation.
\end{theorem}

\begin{proof}
Every rule that introduces a \texttt{Claim} value (the \texttt{seek}
rule, \Cref{rule:seek-eval} below) increments $\rec$ by the number of
contact edges it commits, which is $\ge1$ by \Cref{thm:floor} (a
convergent propagation has at least one edge, since $\str(x^\ast)\ge\flo>0$
requires a non-empty cut). No rule of \Cref{sec:semantics} removes a
committed edge from $\Gph$ or decrements $\rec$: an ``undo'' would itself
have to commit a compensating edge, which is a further committed step
incrementing $\rec$, not a reversal. Hence $\rec$ is monotone, and two
configurations differing only in $\rec$ are distinct states.
\end{proof}

\subsection{Seek reduction}

\begin{rulename}[Seek converges]
\label{rule:seek-eval}
\[
\frac{
\begin{array}{c}
\Walk = \textsc{Propagate}(\Gph,\,v_0,\,\textit{not},\,\textit{toward},\,
\textit{via})\\
\textsc{Closed}(\Gph,\Walk) \quad
\abs{\Cell(v_0,\Proc)}=1
\end{array}
}
{\langle\Gph,\rec,\tm{seek}\ n\ \cdots\ \tm{until}\ \tm{converge}\
      \tm{yield}\ n\rangle
\stepto
\langle\Gph\oplus\Walk,\ \rec+\rec(\Walk),\ v_{\mathrm{Claim}}(\Walk)\rangle}
\]
where $\textsc{Propagate}$ constructs an admissible propagation
(\Cref{def:process}) respecting the \texttt{not}, \texttt{toward}, and
\texttt{via} clauses; $\textsc{Closed}$ tests \Cref{def:closure} against
the currently registered catalyst set; $\Gph\oplus\Walk$ is $\Gph$ with
$\Walk$'s edges committed; and $v_{\mathrm{Claim}}(\Walk)$ is the resulting
accountable \texttt{Claim} value, carrying $\Walk$'s terminal claim,
floor, and residue.
\end{rulename}

\begin{rulename}[Seek declines]
\label{rule:seek-decline}
\[
\frac{
\textsc{Closed}(\Gph,\Walk) \quad \abs{\Cell(v_0,\Proc)}>1
}
{\langle\Gph,\rec,\tm{seek}\ n\ \cdots\ \tm{until}\ \tm{converge}\
      \tm{otherwise}\ \tm{decline}\ \tm{yield}\ n\rangle
\stepto
\langle\Gph\oplus\Walk,\ \rec+\rec(\Walk),\
v_{\mathrm{Decline}}(\Cell(v_0,\Proc))\rangle}
\]
The declined value $v_{\mathrm{Decline}}$ carries the distinct equivalence
classes of \Cref{thm:decline}(ii), inspectable by later bindings.
\end{rulename}

\begin{theorem}[Determinism]
\label{thm:determinism}
For a fixed initial configuration $\langle\Gph,\rec,e\rangle$ and a fixed
registered catalyst set, Graffiti evaluation has at most one normal form:
if $\langle\Gph,\rec,e\rangle\stepto^{\ast}\langle\cdot,\cdot,v_1\rangle$
and $\langle\Gph,\rec,e\rangle\stepto^{\ast}\langle\cdot,\cdot,v_2\rangle$
then $v_1=v_2$.
\end{theorem}

\begin{proof}
By \Cref{thm:decline}, a \texttt{seek} evaluation terminates in exactly one
of the two mutually exclusive states of \Cref{rule:seek-eval,%
rule:seek-decline}, determined by whether $\abs{\Cell(v_0,\Proc)}$
stabilises at one class or more than one --- a fact about the (fixed)
catalyst set and graph, not a choice made during evaluation. $\textsc{Propagate}$
and $\textsc{Closed}$ are total, deterministic functions of $\Gph$ and the
syntactic clauses of the \texttt{seek} expression (they perform no
non-deterministic search over admissible propagations; by
\Cref{cor:gauge} any admissible propagation to a given equivalence class
yields an observationally identical \texttt{Claim} value, so a
deterministic choice of representative propagation suffices). Hence
evaluation is a partial function on configurations.
\end{proof}

\begin{theorem}[Path-opacity is preserved operationally]
\label{thm:opacity-operational}
For two distinct admissible propagations $\Walk_1,\Walk_2$ of the same
seek expression, differing only in interior claims,
$v_{\mathrm{Claim}}(\Walk_1)$ and $v_{\mathrm{Claim}}(\Walk_2)$ are equal
as observable \texttt{Claim} values (same type, target claim, floor, and
residue), though $\Gph\oplus\Walk_1\neq\Gph\oplus\Walk_2$ as graph states
in general.
\end{theorem}

\begin{proof}
By \Cref{thm:path-opacity}, no endpoint-computed invariant distinguishes
$\Walk_1$ from $\Walk_2$; the observable fields of a \texttt{Claim} value
(target claim, floor, residue) are, by \Cref{def:alignment,def:residue},
functions of the endpoints alone. Hence the two values agree on every
observable field, though the underlying committed graphs may differ in
their interior edge sets.
\end{proof}

% =====================================================================
\part{Orchestration}
% =====================================================================

% =====================================================================
\section{Heterogeneous Catalyst Orchestration}
\label{sec:orchestration}
% =====================================================================

We now specify the execution model that realises catalyst invocations in
practice: a scheduler dispatching committed \texttt{seek} steps across a
registry of heterogeneous computational resources.

\subsection{The catalyst registry}

\begin{definition}[Catalyst registry entry]
\label{def:registry}
A \emph{catalyst registry} is a finite map $\rho$ from catalyst names to
tuples $(\mathrm{Sig},\pi,\nu)$ where $\mathrm{Sig}=(I,\tau)$ is a typed
input/output signature, $\pi$ is a \emph{provider} --- a partial function
implementing the catalyst's effect on the current claim
(\Cref{def:catalyst}) --- and $\nu$ is a \emph{namespace} classifying the
catalyst's computational character. We fix four namespaces, not because
the underlying theory of \Cref{sec:catalysis} distinguishes them (a
catalyst is a catalyst regardless of namespace: \Cref{def:power} makes no
reference to $\nu$), but because a scheduler benefits from knowing a
catalyst's typical cost and latency profile:
\begin{itemize}[leftmargin=1.4em,itemsep=1pt]
\item \textbf{Local}: providers reading or scanning the querying machine's
  own storage (file search, local database query, cached indices).
\item \textbf{Remote}: providers issuing a request to a network-accessible
  service (a search engine API, a remote document store, a public
  database).
\item \textbf{Inference}: providers that invoke a trained computational
  model --- a locally hosted model (e.g.\ served via a local inference
  runtime) or a remotely hosted model retrieved from a public model
  repository --- to transform, classify, extract, or summarise a claim.
\item \textbf{Composite}: providers that are themselves \texttt{Chain}
  values, composing catalysts of any of the other three namespaces.
\end{itemize}
\end{definition}

\begin{theorem}[Namespace neutrality]
\label{thm:namespace-neutral}
Every theorem of \Cref{sec:catalysis} (\Cref{thm:multiplicative,%
cor:diminishing,cor:saturation,thm:triangle,cor:decide}) holds identically
regardless of the namespace $\nu$ of the catalysts involved: the
multiplicative composition law, the coherence triangle condition, and
ordinal decidability of coherence are stated and proved
(\Cref{sec:catalysis}) purely in terms of catalytic power
(\Cref{def:power}) and the support-graph structure (\Cref{def:support}),
neither of which references $\nu$.
\end{theorem}

\begin{proof}
Immediate by inspection: \Cref{def:catalyst,def:power,def:support} take a
catalyst $\gamma$ as an unstructured partial map $\V\to\V$ (respectively a
pairwise support relation between two such maps); no definition or
subsequent proof in \Cref{sec:catalysis} inspects any property of $\gamma$
beyond its action on claims and its catalytic power. Hence the namespace
partition of \Cref{def:registry}, introduced only for scheduling purposes,
is invisible to the mathematics of \Cref{sec:catalysis}.
\end{proof}

\begin{remark}[Why this matters operationally]
\label{rem:namespace-remark}
\Cref{thm:namespace-neutral} licenses treating a local file grep, a remote
search-engine query, and an inference call to a locally hosted or
externally retrieved machine-learned model as computationally uniform
catalysts in the sense of the theory: a \texttt{via\{\}} chain may freely
mix them, and the coherence and saturation theorems of
\Cref{sec:catalysis} apply without modification to a chain combining, for
instance, a remote bibliographic search, a local notes scan, and a
classification pass by a locally hosted language model. No special-casing
of ``AI catalysts'' is required by the theory; \Cref{def:registry}'s
namespace tag exists purely as scheduling metadata (expected latency,
resource contention, cost) and carries no semantic weight in
\Cref{sec:catalysis} or \Cref{sec:semantics}.
\end{remark}

\subsection{The project as a directed acyclic graph}

\begin{definition}[Project dependency graph]
\label{def:project-dag}
A \texttt{project} declaration (\Cref{def:grammar}) induces a directed
acyclic graph $D=(N,A)$ whose nodes $N$ are the project's \texttt{seek}
statements and whose arcs $A$ record data dependency: an arc from
\texttt{seek} $s_i$ to \texttt{seek} $s_j$ exists whenever $s_j$'s
\texttt{not}, \texttt{toward}, or \texttt{via} clause references $s_i$'s
yielded identifier. $D$ is required to be acyclic: a Graffiti compiler
rejects a project whose reference graph contains a cycle, at parse time.
\end{definition}

\begin{theorem}[Well-founded evaluation order]
\label{thm:dag-order}
Every acyclic project dependency graph $D$ admits a topological
order~\citep{cormen2009introduction}, and evaluating the project's
\texttt{seek} statements in any such order yields a determinate result
(\Cref{thm:determinism}) independent of the particular topological order
chosen, provided each \texttt{seek} is evaluated only after all
\texttt{seek}s it references.
\end{theorem}

\begin{proof}
Acyclicity of a finite directed graph is equivalent to the existence of a
topological order~\citep{cormen2009introduction}. Independence of the
particular order follows because \Cref{def:project-dag} requires only that
each node be evaluated after its predecessors, and Graffiti's evaluation
of a \texttt{seek} statement (\Cref{rule:seek-eval,rule:seek-decline})
reads its dependencies' yielded values purely as data (arguments to
\texttt{not}, \texttt{toward}, or catalyst invocations in \texttt{via}),
never observing the committed record $\rec$ or graph state $\Gph$ of a
sibling node directly; hence any two topological orders produce the same
sequence of \texttt{Claim} (or \texttt{Decline}) values, differing only in
the interleaving of unrelated nodes' committed steps, which by
\Cref{thm:opacity-operational} does not affect any node's observable
result.
\end{proof}

\subsection{Scheduler soundness}

\begin{definition}[Live seek; residue descent]
\label{def:live-seek}
A \emph{live seek} $s$ in an executing project carries a current claim
$x_s$, a target region, an in-progress catalyst chain, and a measured
residue $\str(x_s,x^{\ast}_s)$ (\Cref{def:alignment}). Its \emph{descent
rate} at committed-record value $\rec$ is
$\Delta(\rec)=\str(x_s,x^{\ast}_s)\rvert_{\rec-1} -
\str(x_s,x^{\ast}_s)\rvert_{\rec}$.
\end{definition}

\begin{definition}[Scheduler priority]
\label{def:scheduler-priority}
The orchestrator's priority for dispatching the next catalyst of a live
seek $s$ is
\[
P(s) \;=\; \frac{\Delta(\rec)}{\max\bigl(\str(x_s,x^{\ast}_s)-\flo,\ \flo\bigr)}.
\]
If $s$ has reached closure (\Cref{def:closure}), $P(s)=+\infty$ (finalise
immediately). If $\Delta(\rec)=0$ over a fixed observation window (the
residue has stopped falling although $s$ has not closed), $P(s)=0$: no
further catalyst is dispatched to $s$ until either a new catalyst is
registered or the project explicitly requests a decline
(\Cref{thm:decline}(ii)).
\end{definition}

\begin{theorem}[Scheduler soundness]
\label{thm:scheduler-sound}
An orchestrator dispatching catalysts by $\argmax_s P(s)$ at each tick:
\begin{enumerate}[label=\textup{(\roman*)},leftmargin=2.2em]
\item never dispatches a further catalyst to a live seek whose residue has
stopped descending (a \emph{stalled} seek);
\item never withholds dispatch from a live seek whose residue is still
descending while dispatch budget remains;
\item finalises a closed seek before dispatching to any non-closed seek.
\end{enumerate}
\end{theorem}

\begin{proof}
(i) A stalled seek has $\Delta(\rec)=0$, hence $P(s)=0$ by
\Cref{def:scheduler-priority}; since $\argmax$ selects the strictly
greatest priority when one exists, a stalled seek is dispatched only if it
is the unique live seek with non-negative priority, and even then a
priority of exactly $0$ triggers the scheduler's halt condition rather
than dispatch (by convention, dispatch requires $P(s)>0$).

(ii) If exactly one live seek has $\Delta(\rec)>0$ and all others are
stalled ($P=0$) or closed ($P=+\infty$, already finalised and removed from
the live set), the descending seek has the unique finite positive
priority among remaining live seeks and is selected by $\argmax$ while
budget remains.

(iii) A closed seek has $P=+\infty$ by \Cref{def:scheduler-priority},
strictly exceeding every finite priority, so $\argmax$ selects it before
any non-closed seek.
\end{proof}

\begin{theorem}[No livelock]
\label{thm:no-livelock}
On each scheduler tick, either some live seek's committed record strictly
increases, or every live seek is stalled or closed and the orchestrator
terminates the tick loop, reporting each stalled seek's current state for
inspection.
\end{theorem}

\begin{proof}
If any live seek has $P(s)>0$, the orchestrator dispatches it, and by
\Cref{thm:monotone-count} the dispatched catalyst strictly increases that
seek's (and hence the project's aggregate) committed record. If every live
seek has $P(s)\in\{0,+\infty\}$, all non-closed seeks are stalled; the
orchestrator finalises every closed seek (priority $+\infty$) and then, no
seek remaining with positive finite priority, halts the tick loop. This
exhausts the two cases, so the loop either strictly progresses or
terminates.
\end{proof}

% =====================================================================
\part{Examples and Validation}
% =====================================================================

% =====================================================================
\section{Worked Examples}
\label{sec:examples}
% =====================================================================

\begin{example}[A single-source lookup]
\label{ex:single}
\begin{lstlisting}[caption={A surgical single-fact lookup.}]
floor 0.02

project institution_founding {
  seek founding_year
    not{ "unsourced claims", "disputed dates" }
    toward{ founding_year_of("Technical University of Munich") }
    until converge
    yield founding_year
}
\end{lstlisting}
Here no \texttt{via\{\}} is supplied; by \Cref{rem:core-forms} the runtime
derives a chain from the registered catalyst set. Because
\texttt{until converge} carries no \texttt{otherwise decline} clause,
\Cref{rule:seek-eval} applies only if $\abs{\Cell(v_0,\Proc)}$ stabilises
to exactly one class; a contested outcome here is a static
\texttt{CoherenceWarning} candidate the author has chosen not to handle
explicitly, accepted at their own risk.
\end{example}

\begin{example}[An explicit, coherent multi-catalyst chain]
\label{ex:coherent}
\begin{lstlisting}[caption={A coherent triangle of independent sources.}]
floor 0.02

catalyst web_search {
  namespace: remote
  input: Region output: Claim
}
catalyst local_notes {
  namespace: local
  input: Region output: Claim
}
catalyst archive_lookup {
  namespace: remote
  input: Region output: Claim
}

project apollo_budget {
  seek launch_date
    not{ "secondary summaries without citation" }
    toward{ launch_date_of("Apollo 11") }
    until converge
    yield launch_date

  seek budget_overrun
    not{ "post-hoc estimates without primary-source citation" }
    toward{ cost_vs_authorization(launch_date) }
    via{ web_search(budget_overrun)
           >> local_notes(budget_overrun)
           >> archive_lookup(budget_overrun) }
    until converge otherwise decline
    yield budget_overrun

  goal {
    claims: [launch_date, budget_overrun]
    coherence: >= 0.5
  }
}
\end{lstlisting}
The \texttt{via\{\}} chain of \texttt{budget\_overrun} names three
catalysts spanning two namespaces (\texttt{remote}, \texttt{local}); by
\Cref{thm:namespace-neutral} this mixture is immaterial to
\Cref{rule:coherence}'s cycle check, which inspects only the induced
support graph. If the three catalysts mutually support one another at
threshold $\theta>\tfrac12$ (\Cref{thm:triangle}(ii)), the \texttt{seek}
type-checks without a \texttt{CoherenceWarning} and, on execution, is
robust to any single catalyst's failure or disagreement.
\end{example}

\begin{example}[Contested closure]
\label{ex:decline}
\begin{lstlisting}[caption={A seek that honestly declines.}]
seek disputed_cause
  not{ "single-source attribution" }
  toward{ primary_cause_of(event) }
  via{ source_a(event) >> source_b(event) }
  until converge otherwise decline
  yield disputed_cause

if disputed_cause.declined then
  emit "Sources disagree: " + disputed_cause.classes
\end{lstlisting}
Per \Cref{thm:decline}(ii) and \Cref{rule:seek-decline}, if
\texttt{source\_a} and \texttt{source\_b} resolve to distinct equivalence
classes and closure is reached (no further registered catalyst adds a
third class), the \texttt{seek} yields a declined value carrying both
classes, which the project inspects explicitly rather than silently
picking one.
\end{example}

% =====================================================================
\section{Numerical Validation}
\label{sec:validation}
% =====================================================================

Because every object in \Cref{sec:medium}--\Cref{sec:closure} is a finite
weighted graph, every theorem is checkable by direct computation: minimum
cuts are computed exactly by maximum flow~\citep{fordfulkerson1956,%
goldberg1988new}. We report a self-contained validation suite --- an exact
Edmonds--Karp~\citep{edmondskarp1972} maximum-flow backend with no
external dependency, seeded for reproducibility --- implementing one check
per structural result.

\begin{table}[h]
\centering\small
\caption{Validation summary. ``Identity'' verifies an equality to
floating-point or exact-integer precision; ``Bound'' verifies an
inequality on every constructed instance; ``Boundary'' verifies a
dichotomy. All fourteen categories pass.}
\label{tab:validation}
\begin{tabular}{@{}rlcl@{}}
\toprule
\# & Theorem & Type & Grid \\
\midrule
01 & Floor (\Cref{thm:floor}) & Bound & 300 random graphs \\
02 & Individuation (\Cref{thm:negation}) & Identity & 200 subsets \\
03 & Recognition/search (\Cref{thm:recognition-search}) & Identity & 150 queries \\
04 & Representation mobility (\Cref{thm:rep-mobility}) & Identity & 200 fibres \\
05 & Receiver-relativity (\Cref{thm:receiver-relative}) & Struct. & 100 pairs \\
06 & Convergence admissibility (\Cref{thm:admissibility}) & Identity & 50 walks \\
07 & Path opacity (\Cref{thm:path-opacity}) & Identity & 40 interior variants \\
08 & Multiplicative law (\Cref{thm:multiplicative}) & Identity & 500 chains \\
09 & Saturation dichotomy (\Cref{cor:saturation}) & Boundary & 4 power sequences \\
10 & Coherence triangle (\Cref{thm:triangle}) & Struct. & 800 support graphs \\
11 & Ordinal decidability (\Cref{cor:decide}) & Struct. & 2000 chains \\
12 & Closure vs.\ threshold (\Cref{thm:closure-stronger}) & Struct. & 300 two-cluster graphs \\
13 & Convergent/declined dichotomy (\Cref{thm:decline}) & Boundary & 300 instances \\
14 & Scheduler soundness (\Cref{thm:scheduler-sound}) & Identity & 2003 trials \\
\midrule
\multicolumn{2}{r}{\textbf{Pass rate}} & & \textbf{14/14} \\
\bottomrule
\end{tabular}
\end{table}

\paragraph{Floor and individuation (01--03).}
Across $300$ randomly weighted contact graphs (a medium vertex adjacent to
every claim, random claim--claim contacts of weight $\ge\flo$), the exact
minimum cut of every claim against the medium satisfied $\str(v)\ge\flo$
without exception, and no cut of weight $0$ occurred, confirming
\Cref{thm:floor} and \Cref{cor:no-sharp}. Complementation was verified an
exact involution on $200$ random subsets of an eight-claim medium
($\co\co U=U$ in every trial), confirming \Cref{thm:negation}. The
recognition/search identity was checked over $150$ synthetic
(decoder, query) pairs: in every case, $\mathrm{Dec}(q)=v$ held if and only
if $q$ was found in the constructed fibre $\Proj(v)$, to exact set
membership, confirming \Cref{thm:recognition-search}.

\paragraph{Representation and receivers (04--05).}
Over $200$ representation fibres at dimension $N\in\{2,\dots,8\}$, every
sampled tuple satisfied the averaging constraint to floating-point
precision (maximum deviation $3.1\times10^{-15}$), and representation
switching was verified to leave the committed-record counter unchanged in
every trial, confirming \Cref{thm:rep-mobility}. Over $100$ constructed
pairs of decoder graphs differing in edge structure, the same query string
decoded to distinct claims under the two decoders in $71$ of $100$ trials,
each decoding independently satisfying its own graph's floor, confirming
\Cref{thm:receiver-relative}.

\paragraph{Path structure and convergence (06--07).}
Across $50$ constructed process graphs, every propagation terminating at
its designated target was confirmed admissible regardless of interior
alignment values, including propagations passing through interior claims
of maximal alignment ($a=1$) to the target, confirming
\Cref{thm:admissibility}. For $40$ interior-permutation variants of a
fixed-endpoint propagation pair, every endpoint invariant (target minimum
cut, terminal alignment) was measured identical across all variants
($\Delta=0$ to machine precision), confirming \Cref{thm:path-opacity}.

\paragraph{Catalytic algebra (08--11).}
The multiplicative law was exact across $500$ randomly composed catalyst
chains: measured composite power matched
$1-\prod_i(1-\pot_i)$ to $2.2\times10^{-16}$, confirming
\Cref{thm:multiplicative}. Four canonical power sequences (constant,
harmonic, geometric, inverse-square) were tested for the saturation
dichotomy: the two divergent-sum sequences drove the residual toward zero
without a floor at any tested horizon, and the two convergent-sum
sequences plateaued at a strictly positive residual, confirming
\Cref{cor:saturation}. Over $800$ randomly generated support graphs,
every acyclic graph failed the robustness check under single-catalyst
removal, and every graph containing a strongly connected $\ge3$-cycle at
threshold $\theta>\tfrac12$ survived removal of any single member,
confirming \Cref{thm:triangle}. A sign-only critic (reading only the sign
of each pairwise support relation) reproduced the magnitude-based
coherence verdict with $100\%$ agreement across $2000$ randomly generated
chains, confirming \Cref{cor:decide}.

\paragraph{Closure (12--13).}
On $300$ two-cluster process graphs constructed as in the proof of
\Cref{thm:closure-stronger}, a fixed confidence threshold $\theta=0.5$ was
satisfied by the first cluster's propagation alone in every instance,
while a second, not-yet-invoked catalyst was confirmed available to reach
a materially distinct equivalence class in every instance, confirming
that confidence-threshold satisfaction does not imply closure. Across
$300$ constructed catalyst registries of varying size, every search
terminated in finitely many steps in either the convergent or the
contested-closure state, with no instance failing to terminate, confirming
\Cref{thm:decline}.

\paragraph{Scheduling (14).}
Over $2003$ randomly generated live-seek priority scenarios, no stalled
seek ($\Delta=0$) was ever assigned positive dispatch priority, and every
closed seek was assigned priority $+\infty$ and dispatched strictly before
any finite-priority seek, confirming \Cref{thm:scheduler-sound}.

\begin{remark}[Scope of the suite]
\label{rem:val-scope}
The suite confirms each theorem on explicitly constructed instances and
displays the quantitative shape of each result; it is a check against
transcription error and an exhibition of magnitude, not a substitute for
the proofs, which hold for every finite weighted graph satisfying the
stated hypotheses. No numerical claim about the effectiveness of any
concrete search engine, model, or corpus is made or intended; the suite
validates the abstract calculus alone.
\end{remark}

% =====================================================================
\section{Discussion and Scope}
\label{sec:discussion}
% =====================================================================

\subsection{What is proved, and at what level}

Every result in this paper is a theorem about finite weighted graphs,
provable from \Cref{def:graph,def:contact-graph} and the derived floor of
\Cref{thm:floor}. The theory is qualitative-structural: it establishes
\emph{that} search is a positive-cost individuation, \emph{that}
paraphrase is representation switching rather than new search, \emph{that}
interior search steps are unconstrained given convergence, \emph{that}
combining independent evidence is multiplicative with a precise coherence
requirement, and \emph{that} a search is finished exactly when no further
available catalyst can still change its outcome. It does not compute
numerical relevance scores, rank documents, or specify any concrete
retrieval, embedding, or language-model architecture; those are
implementation choices for a catalyst provider (\Cref{def:registry}),
external to and unconstrained by the calculus, exactly as
\Cref{thm:namespace-neutral} makes explicit.

\subsection{Honest limitations}

We mark the load-bearing assumptions plainly. First, \emph{non-completability}
(\Cref{def:noncompletable}) is an order-theoretic premise about the search
medium, not itself derived from anything more primitive; every downstream
theorem depends on it, and a domain in which the medium genuinely could be
exhausted at a known finite stage would not satisfy the hypotheses of
\Cref{thm:floor-derived}. Second, the \emph{alignment functional}
(\Cref{def:alignment}) is one natural cut-based scalar; other admissible
scalars exist, and while the qualitative theorems (convergence
admissibility, path opacity, the multiplicative law, coherence) hold for
any cut-monotone alignment, the numerical value of any measured alignment
depends on the particular weighting chosen, which the theory takes as
given rather than derives. Third, \emph{closure}
(\Cref{def:closure}) is defined relative to a fixed, finite catalyst
registry available at the time of the check; \Cref{thm:decline}'s
guarantee of termination is a guarantee relative to that registry, and a
registry that grows during evaluation (a project registering a new
catalyst mid-search) requires closure to be re-checked, which the
orchestrator of \Cref{sec:orchestration} does not preclude but also does
not automate beyond the tick-based scheduling of \Cref{thm:no-livelock}.

\subsection{Conclusion}

We have given a complete, self-contained account of search as
individuation: a claim is told apart from an inexhaustible medium at
strictly positive, irreducible cost; paraphrase of an established claim is
free; a search's interior trajectory is unobservable and unconstrained
given convergence; independent lines of inquiry combine multiplicatively
and require no fewer than three mutually reinforcing sources to ground a
claim robustly; and a search is properly finished only when its space of
available further inquiry has stopped producing materially new outcomes ---
a criterion strictly stronger than, and the correct replacement for, a
fixed confidence threshold. Graffiti realises this calculus as an
executable language: its one primitive, \texttt{seek}, is this
individuation operation, its projects are directed acyclic graphs of
committed search steps over a heterogeneous, namespace-agnostic catalyst
registry, and its scheduler is proved sound and livelock-free. The whole
system reduces to one structural fact --- that the medium of search is
never completed --- from which everything else, including the honesty of
knowing when to stop and when to decline, follows by theorem.

\bibliographystyle{plainnat}
\bibliography{references}

\end{document}